\chapter{Introducción general}
\mtcsettitle{minitoc}{Contenidos} 
\minitoc
\section{Contexto y motivación del estudio}
\label{sec:contexto}

\subsection{Orígenes de la pesca comercial y de la industria naval pesquera en Argentina}

La actividad pesquera comercial en Argentina cuenta con antecedentes documentados
desde 1821, con un primer polo de desarrollo en Puerto Rawson, en la provincia de
Chubut \citep{phantom_pesca_argentina}. Sin embargo, el despegue de la pesca como
industria organizada se produce recién en el siglo~XX, de la mano de la
inmigración europea, principalmente española e italiana, que se asentó en
Mar del Plata en las décadas posteriores a la Segunda Guerra Mundial
\citep{phantom_pesca_argentina}. Aquellos primeros pescadores operaban
embarcaciones artesanales, en algunos casos aún propulsadas a vela, orientadas a
capturas de especies pelágicas de media agua (anchoíta, caballa, bonito) que
alimentaron una incipiente industria conservera. El crecimiento del puerto de
Mar del Plata, inaugurado oficialmente en 1922, y la consolidación de sus
escolleras, ordenaron el espacio necesario para que talleres, varaderos y
astilleros artesanales fueran naciendo alrededor de la actividad
\citep{confluencia_portuaria_mdp}. Familias de inmigrantes, como la de Luis
Solimeno, instalado en la ciudad en 1942, contribuyeron decisivamente a
convertir a Mar del Plata en el principal puerto pesquero del país
\citep{argenports_solimeno}.

El salto tecnológico hacia una flota de mayor porte y capacidad se dio a
mediados de la década de 1960, cuando los empresarios Francisco ``Paco''
Ventura y José Greco introdujeron en el mercado interno el filet fresco de
merluza. En 1965, Greco encargó la construcción de cuatro buques de acero en
astilleros de Buenos Aires, cada uno con capacidad para transportar entre
1500 y 2000 cajones de pescado \citep{phantom_pesca_argentina}, marcando el
inicio de una flota industrial de altura construida en el país. El Estado
argentino acompañó, de manera intermitente, este proceso mediante líneas de
crédito blando para la construcción de buques y la incorporación de
tecnología, desde el Decreto 10032/66 del Banco Industrial hasta,
décadas más tarde, el Decreto 145/2019 sobre Lineamientos para la Renovación
de la Flota \citep{phantom_pesca_argentina}.

La industria pesquera nacional experimentó un fuerte crecimiento durante la
década de 1990, impulsado tanto por el aumento del número de buques como por
su mayor potencial de captura: las capturas marítimas totales pasaron de
545 mil toneladas en 1990 a 1{,}34 millones de toneladas en 1997
\citep{rotoplas_historia_pesca}. Este período estuvo acompañado por la
consolidación de la soberanía nacional sobre los recursos vivos de la Zona
Económica Exclusiva argentina y por políticas de fomento a la radicación de
plantas procesadoras, particularmente en la Patagonia
\citep{pesca_maritima_1989_2013}.

Paralelamente, la industria de construcción naval argentina, de la cual la
construcción de buques pesqueros es hoy uno de sus segmentos más dinámicos,
tiene antecedentes que se remontan a la época colonial: el primero data de
1520, cuando Hernando de Magallanes estableció un taller de reparaciones en
la ría de San Julián, y en 1527 se instaló un astillero en el Fuerte Sancti
Spiritu, sobre el río Paraná \citep{argentina_gob_industria_naval}. Ya en el
siglo~XX, hitos como la creación de los Talleres Navales de La Boca en 1936 y
la instalación de infraestructura de reparación sobre el río Luján
consolidaron un sector que, sin embargo, sufrió una crisis profunda a
comienzos de la década de 1990 con la disolución de la flota estatal, el
cierre de numerosos astilleros y talleres, y la pérdida de competitividad
frente a los astilleros asiáticos \citep{padua_construccion_naval}. La
derogación, hacia fines de esa década, del régimen que permitía el
arrendamiento de buques extranjeros para el cabotaje nacional dio inicio a un
proceso de recuperación que, en el segmento pesquero, se profundizó
notablemente en los últimos años.

\subsection{Estado actual de la industria pesquera y de la industria naval pesquera nacional}

La pesca continúa siendo hoy uno de los complejos exportadores más relevantes
de la economía argentina. Según datos de la Subsecretaría de Pesca y
Acuicultura, en 2025 el sector generó divisas por 2010 millones de dólares, un incremento del 4{,}3\,\% respecto de 2024, sobre un total de 549\,416
toneladas exportadas, constituyendo el segundo mejor registro histórico del
complejo pesquero nacional, solo superado por el récord de 2018
\citep{argentina_gob_pesca_2025}. El langostino patagónico (\textit{Pleoticus
muelleri}) es el principal producto de exportación, con 867 millones de
dólares y 119\,775 toneladas en 2025, a un precio promedio de 7240
dólares por tonelada; le siguen el calamar (\textit{Illex argentinus}), cuyas
exportaciones crecieron un 32{,}5\,\% en valor respecto del año anterior, y la
merluza (\textit{Merluccius hubbsi}), con 326 millones de dólares y 127\,422
toneladas \citep{pescare_exportaciones_2025}. China, España y Estados Unidos
concentraron en conjunto algo más de la mitad de las exportaciones pesqueras
argentinas de ese año \citep{argentina_gob_pesca_2025}.

Este dinamismo exportador, sostenido en gran medida por el crecimiento de la
pesquería de langostino (que pasó de capturas del orden de las 20\,000
toneladas anuales a picos de más de 250\,000 toneladas en apenas un par de
décadas, al punto de ser denominado por actores del sector como ``la soja de
la pesca''), tuvo un correlato directo sobre la industria naval nacional
\citep{argenports_industria_naval}. Astilleros que hasta entonces se dedicaban
exclusivamente a la reparación debieron encarar obras de infraestructura para
comenzar a construir buques nuevos; astilleros del Delta del Paraná y de
Tigre que tradicionalmente construían embarcaciones de turismo o barcazas
incursionaron por primera vez en la construcción de buques pesqueros, y
establecimientos que habían permanecido inactivos durante años fueron
reactivados \citep{argenports_industria_naval}. En la actualidad, el sector de
la industria naval argentina (que incluye, además del segmento pesquero, la
reparación y construcción para otros usos) se compone de más de 300 empresas
y emplea a alrededor de 10\,000 trabajadores en forma directa
\citep{argenports_industria_naval}.

El polo naval pesquero más importante del país se encuentra en el puerto de
Mar del Plata, donde operan, entre otros, el Astillero Naval Federico
Contessi y Cía.\ S.A. (especializado desde hace décadas en la construcción de
buques pesqueros costeros, de altura, congeladores, palangreros y tangoneros
\citep{contessi_historia}), Astilleros Servicios Portuarios Integrados S.A.
(SPI) y TecnoPesca Argentina (TPA), astilleros a los que en 2023 la autoridad
portuaria otorgó concesiones a veinte años para desarrollar nuevas obras de
infraestructura, en el marco del crecimiento sostenido de la demanda de
construcción de embarcaciones nuevas \citep{globalports_concesiones_mdp}. A
ellos se suman astilleros de otras localidades del litoral marítimo y
fluvial, entre los que se destaca el astillero de la firma JPA~S.A., y
emprendimientos más recientes como Astilleros Patagónicos Integrados, que
buscan descentralizar hacia el sur la capacidad de construcción y reparación
naval \citep{pescare_api_patagonia}. Es en este contexto de renovación de
flota (con nuevas unidades de mayor porte y complejidad, construidas
íntegramente en el país) que se enmarca el interés de la presente tesis: los
buques pesqueros analizados son representativos de dicho proceso.

A pesar de esta vitalidad productiva, la predicción de potencia y el diseño
hidrodinámico de estos buques continúan apoyándose, en gran medida, en
metodologías de extrapolación modelo--buque desarrolladas y validadas
históricamente sobre buques mercantes esbeltos. Los pesqueros que hoy se
construyen en el país pertenecen, en cambio, a una tipología de casco de baja
esbeltez, es decir, de relación eslora-manga reducida ($L/B < 4$, donde $L$ es
la eslora y $B$ la manga), cuyas proporciones difieren sustancialmente de las
de aquellos cascos. Esta tipología no es, sin embargo, una particularidad de
la flota argentina: las series sistemáticas clásicas de buques pesqueros,
desarrolladas en distintos países, incluyen cascos de proporciones similares
\citep{yilmaz1999}, lo que evidencia su amplia difusión internacional. Esta
situación motiva la problemática que se desarrolla en la siguiente sección.

\section{Problemática de los efectos de escala en la predicción de potencia}
\label{sec:problematica}

La problemática que aborda esta tesis presenta dos niveles que conviene
distinguir. El primero es general: los efectos de escala son inherentes a
toda predicción de potencia basada en ensayos con modelos, cualquiera sea la
geometría del buque o del propulsor. El segundo es particular: los métodos
tradicionales con los que se corrigen esos efectos pueden no resultar
adecuados para tipologías de casco alejadas de aquellas sobre las que fueron
desarrollados, como es el caso de los buques pesqueros de baja esbeltez.

\subsection{Los efectos de escala como problema general}

La predicción de la potencia de un buque a partir de ensayos con modelos
físicos obliga a extrapolar resultados obtenidos a números de Reynolds muy
inferiores a los de operación del buque. Esta falta de semejanza dinámica
completa entre modelo y buque, presente en todo ensayo con modelos e
independiente de la forma del casco, da origen a los efectos de escala
(Capítulo~\ref{chap:antecedentes}). Estos afectan a las dos grandes
componentes de la predicción de potencia, que definen los dos ejes de esta
tesis: la resistencia al avance, cuya extrapolación descansa sobre el factor
de forma $(1+k)$ (un coeficiente que mide cuánto excede la resistencia viscosa
del casco a la de una placa plana equivalente), y la propulsión, en la que
los resultados de los ensayos de hélices en aguas abiertas deben escalarse a
los números de Reynolds de la hélice del buque.

Precisamente por su carácter general, estas dificultades han sido objeto de
atención sostenida por parte de la comunidad internacional de hidrodinámica
naval. El Comité Especialista en Métodos Combinados CFD/EFD (es decir,
métodos que integran simulaciones numéricas, CFD, y ensayos experimentales,
EFD) de la 29.ª ITTC (International Towing Tank Conference) relevó los problemas de escala que
afectan a la predicción de potencia y los calificó según su impacto sobre las
tendencias de diseño, su impacto sobre el valor absoluto de la potencia y la
frecuencia con que se presentan en la práctica, así como según la posibilidad
de mejorarlos mediante CFD \citep{ITTC_Combined_2021}. A partir de esa
evaluación, el comité propuso a la comunidad concentrar sus esfuerzos en cinco
temas, entre los cuales figuran los dos que constituyen los ejes de la
presente tesis: la determinación numérica del factor de forma y la mejora de
los métodos de escalado de los ensayos de hélices en aguas abiertas.

En el eje de la resistencia al avance, la determinación del factor de forma
obtuvo la calificación más alta de todos los temas evaluados, tanto por su
impacto como por su potencial de mejora mediante CFD. El comité señaló que el
procedimiento experimental habitual para obtenerlo (el método de Prohaska)
puede fallar para algunos tipos de buque, con errores del orden del 5\,\% y de
hasta el 10\,\% en la potencia total, y que puede además distorsionar las
tendencias de diseño, por ejemplo al comparar formas de casco alternativas.
Por esas razones, seleccionó este tema para un estudio colaborativo
específico, cuyos resultados se discuten en el
Capítulo~\ref{chap:antecedentes}.

En el eje de la propulsión, el comité subrayó que el escalado de los
resultados del ensayo de aguas abiertas a otros números de Reynolds es una
parte crucial de la predicción de potencia y que, si bien el procedimiento
ITTC-78 ofrece formulaciones sencillas para ello, los distintos métodos
disponibles conducen a resultados significativamente diferentes. Identificó,
además, dos dificultades de fondo: la presencia de flujo laminar y de
transición laminar--turbulenta sobre las palas a escala de modelo, cuya
simulación numérica sigue siendo una tarea exigente, y la dependencia del
efecto de escala con la geometría de la pala, que afecta especialmente a
diseños no convencionales, como las hélices de área de pala reducida. En su
evaluación, este tema recibió la calificación máxima en impacto sobre las
tendencias de diseño y en frecuencia de aparición, aunque un impacto menor
sobre el valor absoluto de la potencia, con errores estimados de hasta un
3\,\% que pueden conducir a diseños de hélice subóptimos en los tipos de buque
más comunes. El comité concluyó que se requieren nuevos estudios, en los que
el CFD puede formar parte, para eventualmente actualizar el procedimiento de
escalado recomendado por la ITTC \citep{ITTC_Combined_2021}.

\subsection{Limitaciones de los métodos tradicionales en buques de baja esbeltez}

A los efectos de escala comunes a cualquier buque se suma un problema
particular de ciertas tipologías de casco. Los métodos tradicionales de
extrapolación modelo--buque, así como las correlaciones empíricas asociadas
a ellos, fueron desarrollados y calibrados fundamentalmente sobre buques
mercantes esbeltos, y la evidencia que respalda el enfoque combinado EFD--CFD
proviene, hasta el momento, de cascos de referencia de esbeltez moderada
(Capítulo~\ref{chap:antecedentes}). La propia ITTC advierte que, si bien los
errores de estos métodos quedan en promedio bien compensados por las
correlaciones de cada canal de ensayos cuando se trabaja con tipos de buque
habituales, pueden ser mayores para buques que se apartan de esos tipos
\citep{ITTC_Combined_2021}. Su aplicabilidad a cascos de baja relación
eslora-manga, como los de los buques pesqueros, no está por lo tanto
garantizada, y su uso sin una verificación específica puede conducir a
errores significativos en la predicción de potencia.

\subsection{Enfoque y alcance de la tesis}

En correspondencia con estos dos niveles, la tesis persigue un propósito
también doble. En primer lugar, contribuir de manera general al estudio de los
efectos de escala mediante métodos combinados EFD--CFD en los dos ejes
señalados por la ITTC: en el de la resistencia al avance, a través de la
determinación del factor de forma y de su dependencia con la escala, y en el
de la propulsión, a través del escalado de las hélices en aguas abiertas,
tanto en régimen turbulento como en presencia de transición a escala de
modelo. Estos aportes no se restringen a una tipología particular, ya que se
apoyan también en un buque y en hélices de referencia. En segundo lugar,
examinar en particular qué ocurre en el eje de la resistencia al avance con
los buques pesqueros de baja esbeltez, tomando como casos de estudio tres
pesqueros diseñados y construidos en astilleros argentinos, cuyos modelos
fueron ensayados en el canal de remolque del LabHiNO. Si bien los casos de
estudio provienen de la industria naval nacional, la difusión internacional de esta tipología de casco hace que sus resultados
sean de interés para buques pesqueros de proporciones similares en otras
flotas.

\section{Relevancia del estudio}

La predicción de la potencia propulsiva es una etapa central del proyecto de
un buque, ya que a partir de ella se seleccionan, en las etapas tempranas del
diseño, el motor principal y la hélice. Una subestimación de la potencia puede
dar lugar a un buque que no alcance la velocidad de servicio prevista o que
opere con márgenes insuficientes, mientras que una sobrestimación se traduce
en plantas propulsoras sobredimensionadas y en mayores costos de construcción
y de operación a lo largo de la vida útil del buque. La precisión de los
métodos de extrapolación modelo--buque tiene, por lo tanto, consecuencias
técnicas y económicas directas.

En el caso de los buques pesqueros, esa precisión depende de que los métodos
empleados resulten adecuados para cascos de baja esbeltez, cuyas proporciones
difieren de las de los buques sobre los que dichos métodos fueron
desarrollados. Caracterizar los efectos de escala en este tipo de cascos
contribuye, en consecuencia, a una predicción de potencia más confiable en
su diseño, tanto en el ámbito nacional como en el de otras flotas pesqueras
con buques de características similares.

\section{Objetivos generales y específicos}

La tesis se organiza en dos ejes, correspondientes a las dos grandes
componentes de la predicción de potencia, cada uno de ellos con su propio
objetivo general.

\subsection*{Objetivos generales}

\textbf{Eje 1. Resistencia al avance: factor de forma.} Caracterizar,
mediante métodos combinados EFD--CFD, los efectos de escala en la
determinación del factor de forma y su influencia en la predicción de la
resistencia al avance, con particular atención a los buques pesqueros de baja
esbeltez.

\textbf{Eje 2. Propulsión: hélices en aguas abiertas.} Caracterizar,
mediante métodos combinados EFD--CFD, los efectos de escala sobre las
prestaciones de hélices en aguas abiertas, tanto en régimen turbulento como
en presencia de transición laminar--turbulenta a escala de modelo.

\subsection*{Objetivos específicos}

\subsubsection*{Eje 1. Resistencia al avance: factor de forma}

\begin{itemize}
    \item Generar una base de datos experimental propia, obtenida en el canal
    de remolque del LabHiNO, de resistencia al avance para tres buques
    pesqueros de diseño nacional (P1, P2 y P3) y para un buque de referencia
    esbelto, el portacontenedores KCS (\textit{KRISO Container Ship}), uno de
    los cascos estándar de validación en hidrodinámica naval, que sirva de
    base de validación para las simulaciones numéricas.
    \item Determinar el factor de forma $(1+k)$ de dichos buques mediante los
    métodos de Prohaska y de doble cuerpo, contrastando ambos enfoques entre
    sí.
    \item Cuantificar la dependencia del factor de forma con el número de
    Reynolds sobre un rango de escalas geométricas que se extiende desde el
    modelo de ensayo hasta la escala buque, e identificar los mecanismos
    físicos responsables de dicha dependencia.
    \item Evaluar la validez de las hipótesis clásicas de extrapolación
    modelo--buque para cascos de baja esbeltez, a partir de los resultados
    obtenidos para los pesqueros P1, P2 y P3.
    \item Cuantificar el impacto de los efectos de escala identificados sobre
    la predicción de potencia efectiva a escala buque, en comparación con el
    procedimiento estándar de extrapolación de la ITTC.
\end{itemize}

\subsubsection*{Eje 2. Propulsión: hélices en aguas abiertas}

\begin{itemize}
    \item Caracterizar los efectos de escala sobre los coeficientes
    hidrodinámicos de una hélice de referencia de la serie Wageningen~B (una
    familia sistemática y ampliamente documentada de hélices, de uso estándar
    como referencia) en aguas abiertas, desde la escala de modelo hasta la
    escala buque, y compararlos con las predicciones del procedimiento
    estándar ITTC-78.
    \item Evaluar sobre una hélice de diseño convencional, pero no seriada,
    cedida por un proveedor privado, el comportamiento en condiciones clásicas
    de ensayo a escala de modelo y los efectos asociados al régimen de
    transición laminar--turbulenta, que se manifiesta cuando la capa límite
    sobre los álabes opera a números de Reynolds bajos con regiones laminares
    de extensión significativa.
    \item Cuantificar el impacto de los efectos de escala identificados sobre
    la predicción de las prestaciones de la hélice a escala buque, en
    comparación con el procedimiento ITTC-78.
\end{itemize}

\section{Estructura de la tesis}

El resto de la tesis se organiza en cinco capítulos adicionales. El
Capítulo~\ref{chap:antecedentes} desarrolla los fundamentos teóricos de la
similitud dinámica y el estado del arte sobre el factor de forma, su
determinación asistida por simulaciones numéricas y los efectos de escala en
hélices, sentando el marco conceptual sobre el que se apoya el resto del
trabajo. El Capítulo~\ref{chap:Metodologia} describe en detalle las
instalaciones experimentales del LabHiNO, los modelos físicos ensayados, el
procedimiento de extrapolación modelo--buque y la configuración numérica
empleada en OpenFOAM, tanto para las simulaciones de resistencia del casco
(con superficie libre y en configuración de doble cuerpo) como para el
análisis de la hélice en aguas abiertas, además de la estrategia de
verificación y validación adoptada. El
Capítulo~\ref{cap:resistencia_factor_de_forma} desarrolla el Eje~1:
presenta el análisis original de resistencia al avance y factor de forma para el buque pesquero P1 y su
buque de referencia KCS, incluyendo el estudio de los efectos de escala y la
influencia de la proa bulbo, la popa sumergida, el modelo de turbulencia y la
relación $L/B$, extendiendo finalmente el análisis a los pesqueros P2 y P3. El
Capítulo~\ref{chap:helices} desarrolla el Eje~2, correspondiente a la
segunda gran componente de la potencia propulsiva: presenta la validación y el análisis de efectos de escala de la
hélice de referencia Wageningen B4-60 y, sobre el propulsor P206, el estudio
de los efectos del régimen de transición laminar--turbulenta
a escala de modelo. Finalmente, el Capítulo~\ref{cap:conclusiones_generales}
integra los resultados de resistencia y propulsión, discute su impacto sobre
la predicción de potencia y sobre la práctica de diseño y ensayo de buques
pesqueros, resume los aportes principales y propone líneas de trabajo futuro.

% -----------------------------------------------------------------------
% NOTA PARA EL AUTOR: Las siguientes claves de cita fueron construidas a
% partir de las fuentes consultadas para esta sección y deben agregarse
% a la base .bib del proyecto (o reemplazarse por las claves ya
% existentes si alguna de estas fuentes ya estuviera citada):
%   phantom_pesca_argentina, confluencia_portuaria_mdp,
%   argenports_solimeno, rotoplas_historia_pesca,
%   pesca_maritima_1989_2013, argentina_gob_industria_naval,
%   padua_construccion_naval, argentina_gob_pesca_2025,
%   pescare_exportaciones_2025, argenports_industria_naval,
%   contessi_historia, globalports_concesiones_mdp,
%   pescare_api_patagonia
% -----------------------------------------------------------------------